\chapter{Conclusion and Future Scope}
\label{ch:conclusion}

\section{Conclusion}

Pure Transformers become slow and memory-heavy on long inputs because
attention compares every pair of tokens and keeps a growing KV cache.
State-space models such as Mamba-2 avoid both costs with a constant-size, but
lossy, state. This seminar studied \NS, which combines the two: 40 Mamba-2
blocks do most of the sequence processing, 8 attention layers act as anchors
for exact recall, and 40 LatentMoE layers provide capacity. LatentMoE,
shared-weight MTP heads and NVFP4 pretraining complete the design.

LatentMoE moves the experts into a 1024-dimensional latent space, cutting
memory and routing traffic fourfold and paying for more experts, so that 120
billion parameters exist while only about 12 billion are active per token.
Shared-weight multi-token prediction heads give the model its own draft
mechanism for speculative decoding, and pretraining in 4-bit NVFP4 shows that
low-precision training is stable even over 25 trillion tokens.

The model sits in the same accuracy class as the best open Transformers of its
size while running up to 2.2$\times$ and 7.5$\times$ faster and scoring 91.6 on
RULER at one million tokens. The hybrid approach brings new engineering
problems, such as rounding drift in the Mamba state, but the broader lesson is
clear: attention does not have to be everywhere. Efficient large models of the
next generation are likely to be hybrids designed together with their
hardware.

In summary, this seminar has shown how a carefully designed hybrid of
state-space layers, a small number of attention layers and a hardware-aware
Mixture-of-Experts can match strong Transformer models in accuracy while being
considerably cheaper to run on long inputs and outputs. As AI systems move
towards agents that read whole codebases, hold long conversations and call
many tools, the cost of context will matter more and more, and designs such as
\NS\ show a practical way forward.

\section{Future Scope}

\begin{itemize}
  \item \textbf{Larger hybrids.} The Nemotron~3 family includes a larger
        \emph{Ultra} model~\cite{nvidia2025nemotron3}, which will show whether
        the same layout continues to scale.
  \item \textbf{Placement of attention layers.} The ratio of about 8--9\%
        attention comes from experiments on small models; a systematic study
        at large scale, or a learned placement, could improve the trade-off.
  \item \textbf{Quantization-robust state-space layers.} Training methods that
        make the recurrence robust to low-precision storage would remove the
        need for deployment-time fixes and might allow fewer bits for the
        state.
  \item \textbf{Lower-precision training everywhere.} Extending 4-bit training
        to the layers still kept in BF16, and understanding the growth of
        zero-valued gradients, would reduce training cost further.
  \item \textbf{Independent benchmarking.} Reproducing the throughput results
        on other accelerators would show how much of the advantage is
        architectural and how much is specific to one hardware platform.
  \item \textbf{Other modalities and smaller footprints.} Hybrid backbones
        suit long sequences such as audio, video and genomic data, and loading
        experts on demand or pruning rarely used ones could bring such models
        to smaller deployments.
\end{itemize}
